\section{Proofs and Discussion for Section~\ref{sec:theory}}
\label{app:theory}

We first give a closed form for $\varepsilon_\gamma(z)$, then prove
Theorem~\ref{thm:val-par}, and finally prove two further consequences used in
Sections~\ref{sec:theory} and~\ref{sec:evaluation}. Throughout, we use the
closed form of the one-dimensional Wasserstein distance,
$W_1(G,R)=\int_0^1|G^{-1}(u)-F_R^{-1}(u)|\,du$, where $F_R^{-1}$ is the quantile
function of $R$.

\begin{proposition}[Closed form of the compatibility threshold]
\label{prop:compat}
The infimum defining $\varepsilon_\gamma(z)$ in Theorem~\ref{thm:val-par} is
attained, and
\begin{equation}
  \varepsilon_\gamma(z)\;=\;\int_{1-\gamma}^{1}
      \bigl(\tau-F_{R_z}^{-1}(u)\bigr)_+\,du .
  \label{eq:eps_gamma}
\end{equation}
Moreover, $\varepsilon_\gamma(z)=0$ if and only if $R_z([\tau,1])\ge\gamma$.
\end{proposition}

\begin{proof}
If $G([\tau,1])\ge\gamma$, then $G^{-1}(u)\ge\tau$ for all $u>1-\gamma$, so the
integrand of $W_1(G,R_z)$ is at least $(\tau-F_{R_z}^{-1}(u))_+$ on
$(1-\gamma,1]$ and nonnegative elsewhere. Hence $W_1(G,R_z)$ is at least the
right-hand side of~\eqref{eq:eps_gamma}. Equality holds for the nondecreasing
quantile function $G^{-1}(u)=\max\{F_{R_z}^{-1}(u),\tau\}$ on $(1-\gamma,1]$ and
$G^{-1}=F_{R_z}^{-1}$ otherwise, which places mass at least $\gamma$ on
$[\tau,1]$; so the infimum is attained and equals~\eqref{eq:eps_gamma}. The
integrand vanishes almost everywhere on $(1-\gamma,1]$ exactly when
$F_{R_z}^{-1}(u)\ge\tau$ there, i.e., when $R_z([\tau,1])\ge\gamma$.
\end{proof}

\begin{proof}[Proof of Theorem~\ref{thm:val-par}]
If no $\gamma$-valid action exists, $\varepsilon_{\mathcal A}(z,w_0)=+\infty$
and~\eqref{eq:feasibility_bound} holds trivially. Otherwise, for every
$\gamma$-valid action $a$, $P_a$ is a distribution on $[0,1]$ with
$P_a([\tau,1])\ge\gamma$, so it is admissible in the infimum defining
$\varepsilon_\gamma(z)$ and $W_1(P_a,R_z)\ge\varepsilon_\gamma(z)$. Taking the
minimum over the finite set of $\gamma$-valid actions
gives~\eqref{eq:feasibility_bound}. Program~\eqref{eq:program} is feasible if
and only if some $\gamma$-valid action has $W_1(P_a,R_z)\le\varepsilon$, which,
because $\mathcal A(w_0)$ is finite and the minimum is attained, holds if and
only if $\varepsilon\ge\varepsilon_{\mathcal A}(z,w_0)$. Any solution $a^*$ is
feasible and therefore satisfies $W_1(P_{a^*},R_z)\le\varepsilon$. The final
claim is Proposition~\ref{prop:compat}.
\end{proof}

\noindent Because $\varepsilon_\gamma(z)$ depends only on $\hat f$, the
advantaged mediator distribution at $z$, and $(\tau,\gamma)$, and not on the
recipient's factual mediators, the action set, or the cost, cause~(b) in
Section~\ref{sec:theory} can be attributed to the target rather than to the
recourse problem.

\begin{proposition}[Validity alone forces overshoot]
\label{prop:overshoot}
Every $\gamma$-valid plan satisfies $\mathbb E[P_a]\ge\gamma\tau$, so
$D_{\mathrm{post},i}\le\mu_{\mathrm{ref}}(z)-\gamma\tau$. Moreover,
$\varepsilon_\gamma(z)\ge\bigl(\gamma\tau-\mu_{\mathrm{ref}}(z)\bigr)_+$.
\end{proposition}

\begin{proof}
Scores are nonnegative, so
$\mathbb E[P_a]\ge\tau\,P_a([\tau,1])\ge\gamma\tau$. For any $\gamma$-valid $G$,
$W_1(G,R_z)\ge\mathbb E[G]-\mu_{\mathrm{ref}}(z)\ge\gamma\tau-\mu_{\mathrm{ref}}(z)$.
Taking the minimizing $G$ from Proposition~\ref{prop:compat} gives the
second claim.
\end{proof}

\begin{proposition}[What the parity constraint guarantees]
\label{prop:guarantee}
If recipient $i$ is feasible, then
$|D_{\mathrm{post},i}|\le W_1(P_{a^*_i},R_{z_i})\le\varepsilon$, and the same bound
holds for every 1-Lipschitz statistic of the score distribution. If recipient
$i$ abstains, then $D_{\mathrm{post},i}=D^{\mathrm{fact}}_{\mathrm{pre},i}$.
Hence, with feasibility rate $\pi$ among $n$ recipients,
\begin{equation}
  \frac1n\sum_{i=1}^n|D_{\mathrm{post},i}|
  \;\le\;\pi\,\varepsilon+(1-\pi)\,
  \overline{|D^{\mathrm{fact}}_{\mathrm{pre}}|}_{\mathrm{abst}},
  \label{eq:guarantee}
\end{equation}
where the last term averages over abstaining recipients.
\end{proposition}

\begin{proof}
The first claim is Kantorovich--Rubinstein duality applied to the identity map
on $[0,1]$. An abstaining recipient keeps $w_0$, so $P_{a^*_i}$ is a point mass
at $\hat f_0(w_{0,i},z_i)$. Averaging the two cases gives~\eqref{eq:guarantee}.
\end{proof}

\paragraph{Three regimes.}
Proposition~\ref{prop:overshoot} indexes the behaviour of both methods by a
single interpretable quantity, $\mu_{\mathrm{ref}}(z)-\gamma\tau$.
\begin{itemize}[nosep]
  \item \emph{Provably incompatible}, $\mu_{\mathrm{ref}}(z)<\gamma\tau-\varepsilon$:
        every validity-constrained plan, including ordinary recourse, overshoots
        the advantaged reference, and the parity-constrained method must abstain.
        No action set or cost can change this.
  \item \emph{Room for improvement}, $\mu_{\mathrm{ref}}(z)>\gamma\tau+\varepsilon$:
        validity only guarantees $\mathbb E[P_a]\ge\gamma\tau$, so ordinary
        recourse may leave up to $\mu_{\mathrm{ref}}(z)-\gamma\tau$ of the mediated
        disparity unclosed. The parity constraint forces
        $\mathbb E[P_a]\ge\mu_{\mathrm{ref}}(z)-\varepsilon$. This is where the
        parity constraint should yield paired improvements over the baseline.
        Feasibility still depends on $\varepsilon_\gamma(z)$ and on reachability.
  \item \emph{Intermediate}: the outcome is decided by the full shape of
        $R_z$, through $\varepsilon_\gamma(z)$, not by its mean.
\end{itemize}
The compatibility diagnostics in Section~\ref{sec:evaluation} estimate each of
these quantities from the same empirical score distributions used to solve
\eqref{eq:program}.
